\documentclass[12pt]{article}
\usepackage[utf8]{inputenc}
\usepackage[T1]{fontenc}
\usepackage{hyperref}
\title{Gaa Language Resources for K-12 Curriculum Development}
\author{Compiled by Claude Sonnet 5 (claude-sonnet-5) for Andr\'as Kornai}
\date{}
\begin{document}
\maketitle

\section{Introduction}
This dossier concerns the Gaa language of Nigeria (ISO 639-3 ttb, Glottocode gaaa1245). Two web searches found no resources for it. The first search returned only material on Ga of Ghana, a different Kwa language (ISO 639-3 gaa). The second returned only Ga'anda (ISO 639-3 gqa, Glottocode gaan1243), another Nigerian language with a different code. Neither is cited here, to avoid mixing languages.

\section{Primary Educational Resources}
None found.

\section{Linguistic and Reference Materials}
None found.

\section{Religious and Community Resources}
None found.

\section{Digital and Audio Resources}
None found.

\section{Curriculum Development Context}
Nothing in the search results shows Gaa being taught in schools, which matches the reference description. Any curriculum work would have to start with primary documentation, for example through Nigerian linguistic institutions or the Ethnologue and Glottolog entries for ttb. I did not find those in search and cannot vouch for them.
\end{document}